\section{Ricci-flat K\"ahler surfaces}
\label{sec:ricci-flat-surface-rc}

Let $X$ be a connected compact complex surface with a Ricci-flat K\"ahler form $\omega$. Write $h$ for its Hermitian metric on $T^{1,0}X$ and $g$ for the associated real Riemannian metric, so that $\omega(S,T)=g(JS,T)$. Ricci-flatness makes the mean curvature of $h$ zero. It does not make every directional curvature zero. In real dimension four, the remaining curvature acts on a three-dimensional space of two-forms; its invertibility supplies precisely the positive directions needed for the conformal construction.

Use $g$ to identify bivectors and two-forms. For the complex orientation, the Hodge star gives
\[
 \Lambda^2=\Lambda^+\oplus\Lambda^-,\qquad
 \Lambda^\pm=\{\zeta:*\zeta=\pm\zeta\}.
\]
Each summand has real rank three. The first contains $\omega$ and the real $(2,0)+(0,2)$-forms; the second consists of real primitive $(1,1)$-forms. These descriptions follow by writing the star in a unitary real frame.

Let $R^\nabla$ be the curvature of the Levi--Civita connection, with
$R^\nabla(S,T)=\nabla_S\nabla_T-\nabla_T\nabla_S-\nabla_{[S,T]}$, and define the self-adjoint curvature operator $\mathscr R$ by
\[
 \langle\mathscr R(S\wedge T),U\wedge V\rangle
       =g(R^\nabla(S,T)V,U).
\]
K\"ahler curvature has type $(1,1)$, so $\mathscr R$ annihilates the real $(2,0)+(0,2)$-forms. Its value on $\omega$ is the Ricci form and is therefore zero. Thus
\[
 \mathscr R=0\oplus A,\qquad A:\Lambda^-\longrightarrow\Lambda^-.
\]
The endomorphism $A$ is symmetric and trace free: the trace of $\mathscr R$ is half the scalar curvature. It is the anti-self-dual Weyl operator in this Ricci-flat setting. Norms of curvature below are Hilbert--Schmidt norms of these endomorphisms.

\Needspace{12\baselineskip}
\begin{lemma}\label{lem:ricci-flat-rc-invertibility}\index{RC-positivity!Ricci-flat surface}\index{Weyl curvature!anti-self-dual part}
At a point $x$, the metric $h$ is RC-positive if and only if $A_x$ is invertible.
\end{lemma}
\begin{proof}
For a real tangent vector $V$, set $v=\tfrac12(V-\mathrm iJV)$. The Chern and Levi--Civita connections agree on a K\"ahler manifold. Their curvature pairings give
\[
 \alpha_{h,v}=\tfrac12\bigl(\mathscr R(V\wedge JV)\bigr)^\flat.
\]
Indeed, $h(v,w)=\tfrac12(g(V,W)-\mathrm i g(JV,W))$ for $w=\tfrac12(W-\mathrm iJW)$, and $R^\nabla$ commutes with $J$. Its skew-adjointness eliminates the real part of $h(F_hv,v)$ and gives the displayed formula.

The anti-self-dual projection of $V\wedge JV$ is
\[
 \Pi(V)=V\wedge JV-\tfrac12|V|^2\omega.
\]
It is nonzero for $V\ne0$ and fills $\Lambda^-\setminus\{0\}$. To see this, diagonalize a nonzero real primitive $(1,1)$-form in a unitary real frame. It has the form $c(e_1\wedge Je_1-e_2\wedge Je_2)$, with $c>0$ after interchanging the two complex lines. Taking $V=\sqrt{2c}\,e_1$ gives exactly this form as $\Pi(V)$. Also $|\Pi(V)|^2=|V|^4/2$.

Consequently, $\alpha_{h,v}\ne0$ for every nonzero $v$ if and only if $\ker A_x=0$. Since $K_{\omega,h}=0$, \cref{lem:semipositive-mean-rc-criterion} gives the assertion.
\end{proof}

\Needspace{12\baselineskip}
\begin{lemma}\label{lem:ricci-flat-nonzero-determinant}\index{Weyl curvature!determinant}
If $g$ is not flat, the open set $\{x:\det A_x\ne0\}$ is nonempty.
\end{lemma}
\begin{proof}
Suppose that $\det A=0$ everywhere. On the open set where $A\ne0$, its three eigenvalues are $\lambda,-\lambda,0$ with $\lambda>0$, hence are distinct and have smooth local eigenframes. We first record the curvature identity with its normalization.

Choose an oriented eigenframe $\xi_1,\xi_2,\xi_3$ of $\Lambda^-$ with $|\xi_i|^2=2$, and denote its eigenvalues by $\lambda,\mu,\nu$. As skew-adjoint endomorphisms, these forms have the quaternionic multiplication rules. Write
\[
 \nabla\xi_1=c\otimes\xi_2-b\otimes\xi_3,\quad
 \nabla\xi_2=-c\otimes\xi_1+a\otimes\xi_3,\quad
 \nabla\xi_3=b\otimes\xi_1-a\otimes\xi_2.
\]
Here $a,b,c$ are real one-forms. The Ricci identity and the contracted Bianchi identity, using $\operatorname{Ric}(g)=0$, give
\[
 \mathrm da+b\wedge c=\lambda\xi_1,\qquad
 \mathrm d\lambda=(\lambda-\mu)\xi_3c+(\lambda-\nu)\xi_2b,
\]
and their simultaneous cyclic permutations. In the second formula the two-forms act as skew-adjoint endomorphisms on one-forms. Put $D=\mathrm d^*\mathrm d=-\operatorname{div}_g\nabla$ on functions. Applying $\mathrm d^*$ to the second formula yields
\[
 D\lambda=2\lambda^2+4\mu\nu
          +2(\nu-\lambda)|b|^2+2(\mu-\lambda)|c|^2.
\]
For example, the calculation uses
\begin{align*}
 \mathrm d^*(\xi_3c)&=-\langle b,\xi_1c\rangle
       +\langle a,\xi_2c\rangle+\langle a,\xi_3b\rangle-2\nu,\\
 \mathrm d^*(\xi_2b)&=-\langle b,\xi_1c\rangle
       -\langle a,\xi_2c\rangle-\langle a,\xi_3b\rangle-2\mu,
\end{align*}
which follow from the connection and curvature formulas above; substitution and the cyclic formulas for $\mathrm d\mu,\mathrm d\nu$ cancel the mixed pairings.

Now $\lambda+\mu+\nu=0$ and $|A|^2=\lambda^2+\mu^2+\nu^2$. Multiplying the three scalar equations by their eigenvalues and adding gives
\[
 \lambda D\lambda+\mu D\mu+\nu D\nu
 =18\lambda\mu\nu-2\bigl((\mu-\nu)^2|a|^2
                  +(\nu-\lambda)^2|b|^2+(\lambda-\mu)^2|c|^2\bigr).
\]
The connection expression also gives
\[
 |\nabla A|^2=|\mathrm d\lambda|^2+|\mathrm d\mu|^2+|\mathrm d\nu|^2
       +2\bigl((\mu-\nu)^2|a|^2+(\nu-\lambda)^2|b|^2
                                      +(\lambda-\mu)^2|c|^2\bigr).
\]
The product rule therefore gives, with $\Delta=\operatorname{div}_g\nabla=-D$,
\begin{equation}\label{eq:ricci-flat-weyl-norm}
 \Delta|A|^2=2|\nabla A|^2-36\det A=2|\nabla A|^2.
\end{equation}
This identity also holds on the interior of $\{A=0\}$. Both sides are smooth functions, and these two open sets have dense union in $X$, so the identity holds everywhere. Integration over $X$ gives $\nabla A=0$.

If $g$ is not flat, $A$ is consequently nonzero everywhere and has three distinct constant eigenvalues. Its eigenlines are parallel; local unit sections of them are parallel as well. The induced connection on $\Lambda^-$ therefore has zero curvature. Under $\mathfrak{so}(4)=\Lambda^+\oplus\Lambda^-$, the two summands commute, while the action of $\Lambda^-$ on itself is faithful. These facts follow directly from the bases $e_1\wedge e_2\pm e_3\wedge e_4$, $e_1\wedge e_3\mp e_2\wedge e_4$, $e_1\wedge e_4\pm e_2\wedge e_3$. Thus zero induced curvature forces $\mathscr R(S\wedge T)$ to lie in $\Lambda^+$ for every $S,T$. It already lies in $\Lambda^-$, so $\mathscr R=0$, a contradiction.
\end{proof}

\Needspace{14\baselineskip}
\begin{theorem}\label{thm:ricci-flat-surface-rc-positive}\index{RC-positivity!nonclosed tangent metric}
If $X$ has a nonflat Ricci-flat K\"ahler metric, then $T^{1,0}X$ admits an RC-positive Hermitian metric whose associated $(1,1)$-form is nonclosed. It admits neither a uniformly RC-positive Hermitian metric nor a metric with positive definite mean curvature relative to any Hermitian form.
\end{theorem}
\begin{proof}
By \cref{lem:ricci-flat-nonzero-determinant,lem:ricci-flat-rc-invertibility}, the original tangent metric $h$ is RC-positive on a nonempty open set. Its mean curvature is zero. Apply \cref{thm:conformal-rc-globalization} to obtain $h_t=\mathrm e^{-tf}h$ which is RC-positive everywhere.

In that construction one may take the Gauduchon form to be $\omega$ itself and the closed coordinate ball to be a proper subset of $X$. Since $P_\omega f=1$ off that ball, $f$ is nonconstant. The associated form is $\omega_t=\mathrm e^{-tf}\omega$, with
\[
 \mathrm d\omega_t=-t\mathrm e^{-tf}\mathrm df\wedge\omega.
\]
On a real four-dimensional symplectic vector space, $\zeta\mapsto\zeta\wedge\omega$ is injective on one-forms, as a unitary real frame shows. Thus $\mathrm d\omega_t\ne0$ somewhere.

Finally, the induced determinant metric of $h$ is flat because its Chern curvature is the Ricci form of $\omega$, which vanishes. The two obstructions follow from \cref{prop:trivial-determinant-rc-obstruction}.
\end{proof}

The construction stays on the same complex surface: the initial form is K\"ahler and Ricci-flat, while the new form is nonclosed and has directional tangent positivity. The theorem assumes the initial Ricci-flat K\"ahler metric. Producing such a metric on a given surface is a separate existence question.

Combined with \cref{thm:bundle-curvature-complement-convexity}, the RC-positive tangent metric of \cref{thm:ricci-flat-surface-rc-positive} bounds the convexity of the complement of the zero section in the projective completion of the tangent bundle. This bundle admits no metric with positive definite mean curvature, so \cref{cor:positive-mean-complement-convexity} does not apply to it.

\Needspace{12\baselineskip}
\begin{corollary}\label{cor:ricci-flat-tangent-complement-convexity}
\index{RC-positivity!tangent complement}\index{Projective bundle!Ricci-flat tangent complement}
If $X$ has a nonflat Ricci-flat K\"ahler metric, then the complement
\[
 \mathbb P(\mathcal O_X\oplus T^{1,0}X)\setminus\mathbb P(\mathcal O_X)
\]
is $3$-convex.
\end{corollary}
\begin{proof}
Apply \cref{thm:bundle-curvature-complement-convexity}, with $n=r=2$ and $s=1$, to the RC-positive tangent metric of \cref{thm:ricci-flat-surface-rc-positive}.
\end{proof}

